\documentclass[../../book.tex]{subfiles}
\begin{document}
\graphicspath{{../refs/images/}}
\chaptertitle{Frequency Response, Part~1: Bode Plots}{8}

% ===========================================================================
\section{Another Analytical Lens}
\label{sec:intro}
% ===========================================================================

By now we have two analytical views of a dynamic system.  The transfer function
$G(s)$ is our algebraic handle on the input/output relationship.  The time
response---step response in particular---is our window into how the system
behaves in the time domain.  The step response is intuitive: push the input,
watch what happens, and describe what you see with a concrete vocabulary (rise
time, settling time, overshoot).

The \emph{frequency response} is a third view.  Like the transfer function and
the time response, it is an analytical abstraction---a mathematical way of
characterizing what the system does to a particular class of inputs.  That
class is sinusoids.

The frequency domain is less immediately intuitive than the time domain.
We do not usually think about everyday cause and effect in terms of sinusoids.
But the abstraction is powerful.  The Bode plot, which is the standard graphical
representation of the frequency response, is one of the most widely used tools
in controls.  This week we build the vocabulary to read and sketch one.

% ---------------------------------------------------------------------------
\subsection{Sinusoidal steady state}
\label{sec:sinusoidal}
% ---------------------------------------------------------------------------

The premise of frequency response is the behavior of a linear system under a
sinusoidal input, \emph{after all transients have died out}.  In that sinusoidal
steady state, the output is also a sinusoid at the same frequency as the input.
What changes between input and output is (1) the amplitude and (2) the phase
angle.

\begin{figure}[htbp]
  \figplaceholder[5in]{Sinusoidal frequency response: (a) mechanical system with
    sinusoidal force input and position output; (b) transfer function as a
    complex multiplier on phasors; (c) input and output waveforms showing
    amplitude scaling and phase shift. [Nise Fig.~10.2]}
  \caption{A sinusoidal input $f(t) = M_i\cos(\omega t + \phi_i)$ produces a
    sinusoidal steady-state output $x(t) = M_o\cos(\omega t + \phi_o)$ at the
    same frequency.  The system scales the amplitude by $M(\omega) = M_o/M_i$
    and shifts the phase by $\phi(\omega) = \phi_o - \phi_i$.}
  \label{fig:sinusoidal}
\end{figure}

Figure~\ref{fig:sinusoidal} illustrates the idea.  The magnitude frequency
response is the amplitude ratio:
\begin{equation}
  M(\omega) = \frac{M_o(\omega)}{M_i(\omega)}
  \label{eq:mag_response}
\end{equation}
The phase frequency response is the phase difference between output and input:
\begin{equation}
  \phi(\omega) = \phi_o(\omega) - \phi_i(\omega)
  \label{eq:phase_response}
\end{equation}
Both are functions of $\omega$.  Together they constitute the \emph{frequency
response} of the system---a complete description of how the system processes
sinusoids at any frequency.

% ---------------------------------------------------------------------------
\subsection{Experimental measurement}
\label{sec:experimental}
% ---------------------------------------------------------------------------

One practical advantage of the frequency-response view is that it can be
measured directly from a physical system without having a model.  Drive the
system with a sinusoidal input at frequency $\omega$, wait for transients to
decay, then measure the amplitude ratio and phase shift at the output.  Repeat
across a range of frequencies and you have an empirical Bode plot.

This is not just a thought experiment.  In the USV lab we identified surge and
yaw dynamics from step responses, but sinusoidal sweep testing is equally
valid---and in some situations preferable, because it probes the system at
each frequency individually.  The Bode plot is simultaneously an analytical
prediction from a transfer function \emph{and} a directly measurable property
of a physical system.  That dual nature is part of what makes it so useful.

% ===========================================================================
\section{The Frequency Response Function}
\label{sec:G_jomega}
% ===========================================================================

% ---------------------------------------------------------------------------
\subsection{Evaluating \texorpdfstring{$G(j\omega)$}{G(jω)}}
\label{sec:eval}
% ---------------------------------------------------------------------------

We can derive the frequency response from the transfer function analytically.
The derivation follows the forced response of $C(s) = G(s)R(s)$ under a
sinusoidal input.  After separating the forced (sinusoidal) component from the
transient response via partial fractions, the result is compact: substitute
$s = j\omega$ into the transfer function.

\begin{equation}
  \boxed{G(j\omega) = G(s)\big|_{s\,=\,j\omega}}
  \label{eq:G_jomega}
\end{equation}

The complex-valued function $G(j\omega)$ encodes the magnitude and phase
directly:
\begin{align}
  M(\omega) &= \bigl|G(j\omega)\bigr| \label{eq:magnitude} \\
  \phi(\omega) &= \angle G(j\omega) \label{eq:phase}
\end{align}

% ---------------------------------------------------------------------------
\subsection{Example: \texorpdfstring{$G(s) = 1/(s+2)$}{G(s) = 1/(s+2)}}
\label{sec:example_10_1}
% ---------------------------------------------------------------------------

Take $G(s) = 1/(s+2)$.  Substituting $s = j\omega$:
\begin{equation}
  G(j\omega) = \frac{1}{j\omega + 2} = \frac{2 - j\omega}{\omega^2 + 4}
  \label{eq:example_gjw}
\end{equation}
The magnitude and phase follow directly:
\begin{align}
  M(\omega) &= \frac{1}{\sqrt{\omega^2 + 4}} \label{eq:example_mag} \\[4pt]
  \phi(\omega) &= -\tan^{-1}\!\left(\frac{\omega}{2}\right) \label{eq:example_phase}
\end{align}

A few reference values:
\begin{center}
\begin{tabular}{@{}lll@{}}
\toprule
$\omega$ (rad/s) & $M(\omega)$ & $\phi(\omega)$ \\
\midrule
0   & $1/2 = 0.5$     & $0^\circ$   \\
2   & $1/\sqrt{8} \approx 0.354$ & $-45^\circ$ \\
$\infty$ & $\to 0$   & $-90^\circ$ \\
\bottomrule
\end{tabular}
\end{center}

The separate magnitude and phase plots are shown in
Figure~\ref{fig:example_bode}.

\begin{figure}[htbp]
  \figplaceholder[5in]{Separate magnitude and phase Bode plots for
    $G(s)=1/(s+2)$: magnitude rolls off from $-6$~dB at DC and phase shifts
    from $0^\circ$ toward $-90^\circ$. [Nise Fig.~10.4]}
  \caption{Frequency response of $G(s) = 1/(s+2)$.  Upper: magnitude
    $20\log M$ vs.\ frequency (rad/s) on a log scale.  Lower: phase
    $\phi(\omega)$ in degrees.  At the break frequency $\omega = 2$~rad/s
    the magnitude is $-3$~dB and the phase is $-45^\circ$.}
  \label{fig:example_bode}
\end{figure}

% ===========================================================================
\section{The Bode Plot}
\label{sec:bode_plot}
% ===========================================================================

% ---------------------------------------------------------------------------
\subsection{Two scales that make sketching possible}
\label{sec:scales}
% ---------------------------------------------------------------------------

The Bode plot uses two scale choices that turn out to be essential for
sketching:
\begin{itemize}
  \item \textbf{Magnitude in decibels (dB).}  Plot $20\log M$ rather than $M$
    itself.
  \item \textbf{Logarithmic frequency axis.}  Plot $\omega$ on a log scale so
    that each decade ($\times 10$) occupies equal horizontal distance.
\end{itemize}

The power of these choices follows from the structure of a typical transfer
function.  A general $G(s)$ is a ratio of polynomials: its magnitude is a
product of individual terms.  Taking the logarithm turns that product into a
sum:
\begin{equation}
  20\log|G(j\omega)| = 20\log K
    + \sum_k 20\log|(j\omega + z_k)|
    - 20m\log|\omega|
    - \sum_i 20\log|(j\omega + p_i)|
  \label{eq:log_sum}
\end{equation}

If each individual term can be approximated as a \emph{straight line} on the
log-frequency axis, the composite Bode plot is just the algebraic sum of those
straight lines.  This is the core insight: \textbf{products of transfer
function terms become sums of straight-line segments}, which can be added
graphically.  The same additivity holds for phase.

% ---------------------------------------------------------------------------
\subsection{Anatomy of a Bode plot}
\label{sec:anatomy}
% ---------------------------------------------------------------------------

A Bode plot is two graphs sharing a common log-frequency horizontal axis:
\begin{itemize}
  \item \textbf{Upper plot}: magnitude $20\log M$ (dB) vs.\ $\log\omega$.
  \item \textbf{Lower plot}: phase $\phi$ (degrees) vs.\ $\log\omega$.
\end{itemize}

Key landmarks to identify on any Bode plot:
\begin{description}
  \item[Break frequencies] The frequencies where the slope of the asymptote
    changes.  Each real pole or zero contributes one break; each second-order
    conjugate pair contributes one.
  \item[Asymptotic slopes] The slope of the straight-line approximation between
    breaks, in dB/decade.  A real pole contributes $-20$~dB/decade; a real
    zero $+20$~dB/decade; a second-order pole pair $-40$~dB/decade.
  \item[DC gain (low-frequency asymptote)] The magnitude as $\omega \to 0$.
    This sets the starting level of the magnitude plot.
\end{description}

% ===========================================================================
\section{Bode Form: Normalizing for Sketching}
\label{sec:bode_form}
% ===========================================================================

Before sketching, convert the transfer function to \emph{Bode form} by
factoring each pole and zero so that the constant term is 1, not the leading
coefficient.  This normalization gives each building block a DC gain of exactly
1 (0~dB), so individual contributions can be added directly.

For a real pole at $s = -a$, factor out $a$:
\begin{equation}
  \frac{1}{s + a} = \frac{1}{a} \cdot \frac{1}{\dfrac{s}{a} + 1}
  \label{eq:bode_form_pole}
\end{equation}
The $1/a$ factor is a constant gain contribution; $1/(s/a + 1)$ is the
normalized building block with break frequency $\omega = a$.

For a real zero at $s = -a$:
\begin{equation}
  s + a = a\!\left(\frac{s}{a} + 1\right)
  \label{eq:bode_form_zero}
\end{equation}

A complete transfer function in Bode form:
\begin{equation}
  G(s) = \underbrace{\frac{K\,z_1 z_2\cdots z_k}{p_1 p_2\cdots p_n}}_{\displaystyle K_0}
  \cdot
  \frac{\left(\dfrac{s}{z_1}+1\right)\cdots\left(\dfrac{s}{z_k}+1\right)}
       {s^m\!\left(\dfrac{s}{p_1}+1\right)\cdots\left(\dfrac{s}{p_n}+1\right)}
  \label{eq:bode_form_general}
\end{equation}
The constant $K_0$ shifts the magnitude curve up or down by $20\log K_0$~dB
across all frequencies; the normalized factors each start at 0~dB.

% ===========================================================================
\section{The Building Blocks}
\label{sec:blocks}
% ===========================================================================

We now work through each building block type: its exact frequency response,
its asymptotic magnitude approximation, and its phase behavior.
Table~\ref{tab:building_blocks} summarizes the rules; Figures~\ref{fig:blocks_s_1s}
and~\ref{fig:blocks_za_1za} show the normalized Bode plots.

\begin{table}[htbp]
\centering
\caption{Bode building-block summary.  Break frequency is $a$ for first-order
  blocks, $\omega_n$ for second-order.}
\label{tab:building_blocks}
\begin{tabular}{@{}lllll@{}}
\toprule
Block & $G(s)$ & Low-freq mag & Slope above break & Phase range \\
\midrule
Gain          & $K$                     & $20\log K$~dB & 0~dB/dec      & $0^\circ$ \\
Differentiator & $s$                   & ---           & $+20$~dB/dec  & $+90^\circ$ (const.) \\
Integrator    & $1/s$                   & ---           & $-20$~dB/dec  & $-90^\circ$ (const.) \\
Real zero     & $s/a + 1$               & 0~dB          & $+20$~dB/dec  & $0^\circ \to +90^\circ$ \\
Real pole     & $1/(s/a + 1)$           & 0~dB          & $-20$~dB/dec  & $0^\circ \to -90^\circ$ \\
2nd-order zero & $s^2/\omega_n^2 + 2\zeta s/\omega_n + 1$ & 0~dB & $+40$~dB/dec & $0^\circ \to +180^\circ$ \\
2nd-order pole & $1/(s^2/\omega_n^2 + 2\zeta s/\omega_n + 1)$ & 0~dB & $-40$~dB/dec & $0^\circ \to -180^\circ$ \\
\bottomrule
\end{tabular}
\end{table}

% ---------------------------------------------------------------------------
\subsection{Gain: \texorpdfstring{$G(s) = K$}{G(s) = K}}
\label{sec:gain_block}
% ---------------------------------------------------------------------------

A pure gain is frequency-independent.  The magnitude is $20\log K$~dB at all
frequencies and the phase is $0^\circ$ (for $K > 0$).  On the Bode plot the
gain shifts the entire magnitude curve up or down by a fixed amount---it
changes the level, not the shape.

% ---------------------------------------------------------------------------
\subsection{Integrator and differentiator: \texorpdfstring{$G(s) = 1/s$
and $G(s) = s$}{G(s) = 1/s and G(s) = s}}
\label{sec:integrator}
% ---------------------------------------------------------------------------

The integrator $G(s) = 1/s$ has no break frequency.  Substituting $s=j\omega$:
\begin{equation}
  G(j\omega) = \frac{1}{j\omega} = \frac{1}{\omega}\angle{-90^\circ}
  \label{eq:integrator}
\end{equation}
In dB the magnitude is $20\log(1/\omega) = -20\log\omega$, a straight line of
slope $-20$~dB/decade passing through 0~dB at $\omega = 1$~rad/s.  The phase
is a constant $-90^\circ$.

The differentiator $G(s) = s$ is the inverse: slope $+20$~dB/decade through
0~dB at $\omega = 1$~rad/s, constant phase $+90^\circ$.

A pole of order $m$ at the origin, $G(s) = 1/s^m$, contributes a slope of
$-20m$~dB/decade and a constant phase of $-90m^\circ$.

\begin{figure}[htbp]
  \figplaceholder[5in]{Normalized Bode plots for (a)~$G(s)=s$ and
    (b)~$G(s)=1/s$: straight-line magnitude and constant phase.
    [Nise Fig.~10.9a,b]}
  \caption{Differentiator ($s$) and integrator ($1/s$): magnitude and phase.
    Both are straight lines at all frequencies with no break.}
  \label{fig:blocks_s_1s}
\end{figure}

% ---------------------------------------------------------------------------
\subsection{Real zero: \texorpdfstring{$G(s) = s/a + 1$}{G(s) = s/a + 1}}
\label{sec:real_zero}
% ---------------------------------------------------------------------------

The normalized real zero has break frequency $a$.  Its frequency response:
\begin{equation}
  G(j\omega) = j\frac{\omega}{a} + 1
  \label{eq:real_zero_gjw}
\end{equation}

\textbf{Magnitude asymptotes.}
\begin{itemize}
  \item \emph{Low frequency} ($\omega \ll a$): $G(j\omega) \approx 1$,
    so $20\log M \approx 0$~dB.
  \item \emph{High frequency} ($\omega \gg a$): $G(j\omega) \approx j\omega/a$,
    so $20\log M \approx 20\log(\omega/a)$---slope $+20$~dB/decade starting at
    $\omega = a$.
\end{itemize}
The two asymptotes meet at $\omega = a$.  The exact magnitude there is
$20\log\sqrt{2} \approx +3$~dB above the asymptote.

\textbf{Phase.}  The phase rises from $0^\circ$ at low frequencies to
$+90^\circ$ at high frequencies, passing through $+45^\circ$ at the break.
The asymptotic approximation draws a straight-line ramp of slope
$+45^\circ$/decade from one decade below the break ($0.1a$) to one decade
above ($10a$), where it levels off.

% ---------------------------------------------------------------------------
\subsection{Real pole: \texorpdfstring{$G(s) = 1/(s/a + 1)$}{G(s) = 1/(s/a + 1)}}
\label{sec:real_pole}
% ---------------------------------------------------------------------------

The normalized real pole is the mirror image of the real zero:
\begin{itemize}
  \item \textbf{Magnitude}: flat at 0~dB for $\omega < a$, rolls off at
    $-20$~dB/decade for $\omega > a$.  The exact magnitude at the break is
    $-3$~dB.
  \item \textbf{Phase}: $0^\circ$ at DC, $-45^\circ$ at the break, $-90^\circ$
    at high frequencies.  Asymptotic approximation: $-45^\circ$/decade ramp
    from $0.1a$ to $10a$.
\end{itemize}

\begin{figure}[htbp]
  \figplaceholder[5in]{Normalized Bode plots for (c)~$G(s)=(s/a+1)$ and
    (d)~$G(s)=1/(s/a+1)$: asymptotic (straight-line) and exact responses.
    [Nise Fig.~10.9c,d]}
  \caption{Real zero ($s/a + 1$) and real pole $1/(s/a + 1)$: normalized
    asymptotic and exact Bode plots.  The maximum deviation between the
    asymptote and the exact curve is 3.01~dB in magnitude and $5.7^\circ$
    in phase, both at the break frequency.}
  \label{fig:blocks_za_1za}
\end{figure}

% ---------------------------------------------------------------------------
\subsection{Second-order factors}
\label{sec:second_order}
% ---------------------------------------------------------------------------

A second-order denominator arises from a conjugate pole pair.  The normalized
form is:
\begin{equation}
  G(s) = \frac{1}{\dfrac{s^2}{\omega_n^2} + 2\zeta\dfrac{s}{\omega_n} + 1}
  \label{eq:second_order_block}
\end{equation}
The break frequency is $\omega_n$; $\zeta$ is the damping ratio.

\textbf{Magnitude asymptotes.}
\begin{itemize}
  \item \emph{Low frequency}: $G(j\omega) \approx 1$, so $20\log M \approx 0$~dB.
  \item \emph{High frequency}: $|G(j\omega)| \approx \omega_n^2/\omega^2$,
    so $20\log M \approx -40\log(\omega/\omega_n)$---slope $-40$~dB/decade
    starting at $\omega_n$.
\end{itemize}

\textbf{Phase.}  The phase falls from $0^\circ$ at DC to $-180^\circ$ at
high frequencies, passing through $-90^\circ$ at $\omega = \omega_n$.
The asymptotic approximation uses a slope of $-90^\circ$/decade from
$0.1\omega_n$ to $10\omega_n$.

\begin{figure}[htbp]
  \figplaceholder[5in]{Bode asymptotes for the normalized second-order factor
    $1/(s^2/\omega_n^2 + 2\zeta s/\omega_n + 1)$: (a)~magnitude; (b)~phase.
    [Nise Fig.~10.13]}
  \caption{Second-order pole factor: asymptotic magnitude (left) and phase
    (right).  The magnitude breaks at $-40$~dB/decade at the natural frequency
    $\omega_n$; the phase swings from $0^\circ$ to $-180^\circ$ through
    $-90^\circ$ at $\omega_n$.}
  \label{fig:second_order_bode}
\end{figure}

\textbf{A note on accuracy.}  Unlike the first-order case---where the asymptote
and the exact curve differ by at most 3~dB in magnitude---the second-order
asymptote can be far off near the break frequency.  For small $\zeta$ the
actual magnitude \emph{peaks} significantly above 0~dB near $\omega_n$ (the
resonance peak).  The asymptote does not capture this.  Use the asymptotic
sketch to understand the shape and locate break frequencies, then use MATLAB
for the exact plot.

% ===========================================================================
\section{Combining Building Blocks: A Worked Example}
\label{sec:example}
% ===========================================================================

We now sketch the Bode plot for a transfer function with a pole at the origin,
two finite poles, and a zero:
\begin{equation}
  G(s) = \frac{K(s+3)}{s(s+1)(s+2)}
  \label{eq:example_tf}
\end{equation}

\subsection*{Step 1: Convert to Bode form}

Factor out the constants from each pole and zero to normalize them to unit
DC gain:
\begin{equation}
  G(s) = K \cdot \frac{3\!\left(\dfrac{s}{3}+1\right)}
                      {s \cdot 1\cdot\left(\dfrac{s}{1}+1\right)
                              \cdot 2\cdot\left(\dfrac{s}{2}+1\right)}
       = \underbrace{\frac{3K}{2}}_{\displaystyle K_0}
         \cdot
         \frac{\left(\dfrac{s}{3}+1\right)}
              {s\!\left(\dfrac{s}{1}+1\right)\!\left(\dfrac{s}{2}+1\right)}
  \label{eq:example_bode_form}
\end{equation}
The DC gain factor $K_0 = 3K/2$ shifts the magnitude by $20\log(3K/2)$~dB.
For $K = 1$, $K_0 = 1.5$ and the shift is $+3.5$~dB.

\subsection*{Step 2: Identify break frequencies}

Reading from the normalized form: breaks at $\omega = 1$ (pole), $\omega = 2$
(pole), and $\omega = 3$ (zero).

\subsection*{Step 3: Determine the starting point}

The pole at the origin gives an initial slope of $-20$~dB/decade.  Evaluate
the low-frequency asymptote at a convenient starting point---one decade below
the lowest break, $\omega = 0.1$~rad/s.  At that frequency only the $1/s$ and
$K_0$ terms are active (all $(s/a+1)$ terms are $\approx 1$):
\begin{equation}
  20\log|G(j0.1)| \approx 20\log\!\left(\frac{K_0}{0.1}\right)
    = 20\log(15) \approx 23.5~\text{dB} \quad (K=1)
  \label{eq:example_start}
\end{equation}

\subsection*{Step 4: Sketch the magnitude using slope contributions}

Draw the initial $-20$~dB/decade line from the starting point, then update
the slope at each break frequency.  Table~\ref{tab:example_mag_slopes} tracks
the running slope.

\begin{table}[htbp]
\centering
\caption{Magnitude slope contributions for $G(s) = K(s+3)/[s(s+1)(s+2)]$.}
\label{tab:example_mag_slopes}
\begin{tabular}{@{}llr@{}}
\toprule
Frequency range & Event & Running slope (dB/dec) \\
\midrule
Start to $\omega = 1$  & Pole at origin (always active) & $-20$ \\
$\omega = 1$ to $2$    & Pole at $-1$ adds $-20$        & $-40$ \\
$\omega = 2$ to $3$    & Pole at $-2$ adds $-20$        & $-60$ \\
$\omega = 3$ onward    & Zero at $-3$ adds $+20$        & $-40$ \\
\bottomrule
\end{tabular}
\end{table}

\begin{figure}[htbp]
  \figplaceholder[5.5in]{Bode magnitude plot for $G(s)=K(s+3)/[s(s+1)(s+2)]$,
    $K=1$: (a)~individual component asymptotes; (b)~composite asymptote.
    [Nise Fig.~10.11]}
  \caption{Magnitude Bode plot for the worked example.  Left: each building
    block drawn separately.  Right: the composite asymptote, obtained by
    summing the slopes.}
  \label{fig:example_magnitude}
\end{figure}

\subsection*{Step 5: Sketch the phase using the $45^\circ$/decade rule}

Each pole contributes a $-45^\circ$/decade phase ramp beginning one decade
below its break and ending one decade above.  The zero contributes
$+45^\circ$/decade over the same span.  The pole at the origin contributes a
constant $-90^\circ$ offset at all frequencies.

Table~\ref{tab:example_phase_slopes} shows the ramp start/stop frequencies
and slope contributions for each term.

\begin{table}[htbp]
\centering
\caption{Phase slope contributions for $G(s) = K(s+3)/[s(s+1)(s+2)]$.
  Each ramp applies over the one-decade window around its break.}
\label{tab:example_phase_slopes}
\begin{tabular}{@{}llcc@{}}
\toprule
Source & Ramp active (rad/s) & Slope (deg/dec) \\
\midrule
Pole at origin & (all $\omega$, constant offset) & $-90^\circ$ offset \\
Pole at $\omega = 1$ & $0.1$ to $10$   & $-45$ \\
Pole at $\omega = 2$ & $0.2$ to $20$   & $-45$ \\
Zero at $\omega = 3$ & $0.3$ to $30$   & $+45$ \\
\bottomrule
\end{tabular}
\end{table}

Sum the contributions at each frequency to build the composite phase portrait.
Figure~\ref{fig:example_phase} shows the result.

\begin{figure}[htbp]
  \figplaceholder[5.5in]{Bode phase plot for $G(s)=K(s+3)/[s(s+1)(s+2)]$:
    (a)~individual component ramps starting from the $-90^\circ$ offset;
    (b)~composite phase. [Nise Fig.~10.12]}
  \caption{Phase Bode plot for the worked example.  The composite phase starts
    at $-90^\circ$ (from the pole at the origin) and is shaped by the ramp
    contributions of the two finite poles and the zero.}
  \label{fig:example_phase}
\end{figure}

% ===========================================================================
\section{Checking Your Sketch with MATLAB}
\label{sec:matlab}
% ===========================================================================

The companion script \texttt{freq\_response\_companion.m} works through all
the examples in this chapter using the Control System Toolbox.  It also uses
the \texttt{sketchbode} utility, which overlays the straight-line asymptotic
approximation on the exact MATLAB Bode plot so you can see directly where the
asymptote is accurate and where it deviates.

\texttt{sketchbode} was written by a student at Olin College in 2007 and has
been part of this course since.  The workflow is:
\begin{enumerate}
  \item Sketch the Bode plot by hand, identifying break frequencies and slopes.
  \item Call \texttt{sketchbode(sys)} to check your break frequencies and
    slopes against the computed asymptote.
  \item Call \texttt{bode(sys)} (or inspect the \texttt{sketchbode} overlay)
    to see the exact frequency response.
\end{enumerate}

For second-order factors with small $\zeta$, step~3 matters more than for
first-order systems: the asymptote near $\omega_n$ can be misleading, and
only the exact plot shows the resonance peak.

\end{document}
